% Estilo académico reutilizable
% Uso: colocar % Estilo académico reutilizable
% Uso: colocar % Estilo académico reutilizable
% Uso: colocar % Estilo académico reutilizable
% Uso: colocar \input{../estilos/estilo_academico} inmediatamente después de \documentclass

% Codificación e idioma
\usepackage{iftex}
\ifPDFTeX
\usepackage[utf8]{inputenc}
\usepackage[T1]{fontenc}
\usepackage{lmodern}
\else
\usepackage{fontspec}
\setmainfont{Latin Modern Roman}
\fi
\usepackage[spanish]{babel}

% Paquetes generales
\usepackage{pgfplots}
\usepackage[a4paper, margin=2.5cm]{geometry}
\usepackage{setspace}
\usepackage{multirow}
\usepackage{tikz}
\usepackage{graphicx}
\usepackage{xcolor}
\usepackage{enumitem}

% Matemática
\usepackage{amsmath, amssymb, amsthm, mathtools}

% Títulos y enlaces
\usepackage{titlesec}
\usepackage[hidelinks]{hyperref}

% Paleta académica
\definecolor{academicblue}{RGB}{0,70,140}

% Compatibilidad: evita error si el motor no define \DeclareUnicodeCharacter
\providecommand{\DeclareUnicodeCharacter}[2]{}
% Fuerza el mapeo correcto de Unicode para signos invertidos en compilaciones PDFLaTeX
\DeclareUnicodeCharacter{00BF}{\textquestiondown}
\DeclareUnicodeCharacter{2013}{--}
\DeclareUnicodeCharacter{2014}{---}
\DeclareUnicodeCharacter{2018}{`}
\DeclareUnicodeCharacter{2019}{'}
\DeclareUnicodeCharacter{201C}{``}
\DeclareUnicodeCharacter{201D}{''}

% Ajustes tipográficos
\onehalfspacing
\setlength{\parskip}{0.6em}
\setlength{\parindent}{0pt}

% Jerarquía visual sobria
\titleformat{\section}{\large\bfseries\color{academicblue}}{\thesection}{0.75em}{}
\titleformat{\subsection}{\normalsize\bfseries\color{academicblue}}{\thesubsection}{0.75em}{}
\titleformat{\subsubsection}{\normalsize\itshape}{\thesubsubsection}{0.75em}{}

% Viñetas redondas en azul de sección
\setlist[itemize]{label=\textcolor{academicblue}{\textbullet}}
 inmediatamente después de \documentclass

% Codificación e idioma
\usepackage{iftex}
\ifPDFTeX
\usepackage[utf8]{inputenc}
\usepackage[T1]{fontenc}
\usepackage{lmodern}
\else
\usepackage{fontspec}
\setmainfont{Latin Modern Roman}
\fi
\usepackage[spanish]{babel}

% Paquetes generales
\usepackage{pgfplots}
\usepackage[a4paper, margin=2.5cm]{geometry}
\usepackage{setspace}
\usepackage{multirow}
\usepackage{tikz}
\usepackage{graphicx}
\usepackage{xcolor}
\usepackage{enumitem}

% Matemática
\usepackage{amsmath, amssymb, amsthm, mathtools}

% Títulos y enlaces
\usepackage{titlesec}
\usepackage[hidelinks]{hyperref}

% Paleta académica
\definecolor{academicblue}{RGB}{0,70,140}

% Compatibilidad: evita error si el motor no define \DeclareUnicodeCharacter
\providecommand{\DeclareUnicodeCharacter}[2]{}
% Fuerza el mapeo correcto de Unicode para signos invertidos en compilaciones PDFLaTeX
\DeclareUnicodeCharacter{00BF}{\textquestiondown}
\DeclareUnicodeCharacter{2013}{--}
\DeclareUnicodeCharacter{2014}{---}
\DeclareUnicodeCharacter{2018}{`}
\DeclareUnicodeCharacter{2019}{'}
\DeclareUnicodeCharacter{201C}{``}
\DeclareUnicodeCharacter{201D}{''}

% Ajustes tipográficos
\onehalfspacing
\setlength{\parskip}{0.6em}
\setlength{\parindent}{0pt}

% Jerarquía visual sobria
\titleformat{\section}{\large\bfseries\color{academicblue}}{\thesection}{0.75em}{}
\titleformat{\subsection}{\normalsize\bfseries\color{academicblue}}{\thesubsection}{0.75em}{}
\titleformat{\subsubsection}{\normalsize\itshape}{\thesubsubsection}{0.75em}{}

% Viñetas redondas en azul de sección
\setlist[itemize]{label=\textcolor{academicblue}{\textbullet}}
 inmediatamente después de \documentclass

% Codificación e idioma
\usepackage{iftex}
\ifPDFTeX
\usepackage[utf8]{inputenc}
\usepackage[T1]{fontenc}
\usepackage{lmodern}
\else
\usepackage{fontspec}
\setmainfont{Latin Modern Roman}
\fi
\usepackage[spanish]{babel}

% Paquetes generales
\usepackage{pgfplots}
\usepackage[a4paper, margin=2.5cm]{geometry}
\usepackage{setspace}
\usepackage{multirow}
\usepackage{tikz}
\usepackage{graphicx}
\usepackage{xcolor}
\usepackage{enumitem}

% Matemática
\usepackage{amsmath, amssymb, amsthm, mathtools}

% Títulos y enlaces
\usepackage{titlesec}
\usepackage[hidelinks]{hyperref}

% Paleta académica
\definecolor{academicblue}{RGB}{0,70,140}

% Compatibilidad: evita error si el motor no define \DeclareUnicodeCharacter
\providecommand{\DeclareUnicodeCharacter}[2]{}
% Fuerza el mapeo correcto de Unicode para signos invertidos en compilaciones PDFLaTeX
\DeclareUnicodeCharacter{00BF}{\textquestiondown}
\DeclareUnicodeCharacter{2013}{--}
\DeclareUnicodeCharacter{2014}{---}
\DeclareUnicodeCharacter{2018}{`}
\DeclareUnicodeCharacter{2019}{'}
\DeclareUnicodeCharacter{201C}{``}
\DeclareUnicodeCharacter{201D}{''}

% Ajustes tipográficos
\onehalfspacing
\setlength{\parskip}{0.6em}
\setlength{\parindent}{0pt}

% Jerarquía visual sobria
\titleformat{\section}{\large\bfseries\color{academicblue}}{\thesection}{0.75em}{}
\titleformat{\subsection}{\normalsize\bfseries\color{academicblue}}{\thesubsection}{0.75em}{}
\titleformat{\subsubsection}{\normalsize\itshape}{\thesubsubsection}{0.75em}{}

% Viñetas redondas en azul de sección
\setlist[itemize]{label=\textcolor{academicblue}{\textbullet}}
 inmediatamente después de \documentclass

% Codificación e idioma
\usepackage{iftex}
\ifPDFTeX
\usepackage[utf8]{inputenc}
\usepackage[T1]{fontenc}
\usepackage{lmodern}
\else
\usepackage{fontspec}
\setmainfont{Latin Modern Roman}
\fi
\usepackage[spanish]{babel}

% Paquetes generales
\usepackage{pgfplots}
\usepackage[a4paper, margin=2.5cm]{geometry}
\usepackage{setspace}
\usepackage{multirow}
\usepackage{tikz}
\usepackage{graphicx}
\usepackage{xcolor}
\usepackage{enumitem}

% Matemática
\usepackage{amsmath, amssymb, amsthm, mathtools}

% Títulos y enlaces
\usepackage{titlesec}
\usepackage[hidelinks]{hyperref}

% Paleta académica
\definecolor{academicblue}{RGB}{0,70,140}

% Compatibilidad: evita error si el motor no define \DeclareUnicodeCharacter
\providecommand{\DeclareUnicodeCharacter}[2]{}
% Fuerza el mapeo correcto de Unicode para signos invertidos en compilaciones PDFLaTeX
\DeclareUnicodeCharacter{00BF}{\textquestiondown}
\DeclareUnicodeCharacter{2013}{--}
\DeclareUnicodeCharacter{2014}{---}
\DeclareUnicodeCharacter{2018}{`}
\DeclareUnicodeCharacter{2019}{'}
\DeclareUnicodeCharacter{201C}{``}
\DeclareUnicodeCharacter{201D}{''}

% Ajustes tipográficos
\onehalfspacing
\setlength{\parskip}{0.6em}
\setlength{\parindent}{0pt}

% Jerarquía visual sobria
\titleformat{\section}{\large\bfseries\color{academicblue}}{\thesection}{0.75em}{}
\titleformat{\subsection}{\normalsize\bfseries\color{academicblue}}{\thesubsection}{0.75em}{}
\titleformat{\subsubsection}{\normalsize\itshape}{\thesubsubsection}{0.75em}{}

% Viñetas redondas en azul de sección
\setlist[itemize]{label=\textcolor{academicblue}{\textbullet}}
